\section{Preliminaries}\label{sec:prelim}

We establish notation and recall the background results used throughout
the paper.

\subsection{Notation}\label{ssec:notation}

We write $I$ or $I_d$ for the $d \times d$ identity matrix,
$\GL_d(\R)$ for the general linear group, and $\gl_d(\R)$ for its Lie
algebra of all $d \times d$ real matrices.  For $A, B \in \gl_d(\R)$,
the \emph{commutator} (Lie bracket) is $\comm{A}{B} = AB - BA$.  The
Frobenius norm is $\frobnorm{A} = (\Tr(A^\top A))^{1/2}$, and the
operator norm is $\opnorm{A} = \sup_{\opnorm{x}=1} \opnorm{Ax}$.  We
write $\E[\cdot]$ for expectation and $\sigma^2$ for variance.  The
notation $f = \Theta(g)$ means $c_1 \abs{g} \le \abs{f} \le c_2 \abs{g}$
for positive constants $c_1, c_2$.

For a finite group $G$ and a probability measure $\mu$ on $G$, the
\emph{lazy random walk} driven by $\mu$ has transition operator
$P = \tfrac{1}{2} I + \tfrac{1}{2} T_\mu$, where $(T_\mu f)(g)
= \sum_{h} \mu(h) f(gh)$.

\subsection{Proof state spaces and error operators}

\begin{definition}\label{def:proof_state}
A \emph{proof state space} is a pair $(S, d)$ where $S \subseteq \R^d$
is a subset containing the origin and $d$ is the Euclidean metric.
The \emph{correct proof state} is $s^* = 0$.
\end{definition}

\begin{definition}\label{def:error_op}
An \emph{error operator} is a matrix $E = I + \eps A \in \GL_d(\R)$
where $A \in \gl_d(\R)$ is the \emph{error direction} and $\eps > 0$
is the \emph{error magnitude}.  Given a collection of error operators
$\{E_k\}_{k \ge 1}$, the \emph{error group} is
$G = \langle \{E_k\} \rangle \le \GL_d(\R)$.
\end{definition}

\begin{definition}\label{def:error_prop}
An \emph{$n$-step error propagation} is the cumulative product
\[
  E^{(n)} = \prod_{k=1}^{n} E_k = E_n E_{n-1} \cdots E_1.
\]
The \emph{Lyapunov exponent} of the sequence is
\[
  \lambda = \lim_{n \to \infty} \frac{1}{n}\, \E\bigl[\log \opnorm{E^{(n)}}\bigr],
\]
provided the limit exists.
\end{definition}

\begin{definition}\label{def:abelian}
The error propagation is \emph{abelian} if all error operators commute:
$E_i E_j = E_j E_i$ for all $i, j$.  Equivalently, the error directions
$\{A_k\}$ are \emph{simultaneously diagonalizable} over $\R$
(or over $\mathbb{C}$ if we allow complex eigenvalues).
\end{definition}

\begin{definition}\label{def:spectral_gap}
For a random walk on a finite group $G$ with transition operator $P$
acting on $L^2(G)$, the \emph{spectral gap} is
\[
  \gamma = 1 - \rho\bigl(P\big|_{\mathbf{1}^\perp}\bigr),
\]
where $\rho(\cdot)$ denotes the spectral radius and $\mathbf{1}^\perp$
is the orthogonal complement of the constant functions.
\end{definition}

\subsection{Background results}

We collect several known results that are used in our proofs.

\begin{theorem}[Barrau--Bonnabel {\cite[Theorem~2]{barrau2014}}]
\label{thm:bb}
Let $\chi_t$ be a trajectory on a matrix Lie group $G$ satisfying
$\dot{\chi}_t = f_{u_t}(\chi_t)$.  If the dynamics satisfy the group
affine property
\[
  f_u(ab) = f_u(a)\,b + a\,f_u(b) - a\,f_u(\mathrm{Id})\,b
  \quad \text{for all } u, a, b \in G,
\]
then the Lie logarithm $\xi_t = \log(\eta_t)$ of the invariant error
$\eta_t$ satisfies the linear ODE $\dot{\xi}_t = A_t \xi_t$, with
$\eta_t = \exp(\xi_t)$ holding for arbitrarily large errors.
\end{theorem}

\begin{theorem}[Furstenberg {\cite{furstenberg1963}}]\label{thm:furstenberg}
Let $\{M_k\}_{k \ge 1}$ be i.i.d.\ random matrices in $\GL_d(\R)$
with $\E[\log^+ \opnorm{M_1}] < \infty$.  If the closed group generated
by the support of the distribution of $M_1$ is non-compact and acts
strongly irreducibly on $\R^d$, then the top Lyapunov exponent
$\lambda_1 = \lim_{n \to \infty} \frac{1}{n} \E[\log\opnorm{M_n \cdots M_1}]$
is strictly positive.
\end{theorem}

\begin{theorem}[Diaconis upper bound lemma {\cite[Lemma~2]{diaconis1988}}]
\label{thm:diaconis_lemma}
Let $\mu^{*k}$ denote the $k$-fold convolution of a probability measure
$\mu$ on a finite group $G$ and let $U$ be the uniform distribution.
Then
\[
  \bigl\lVert \mu^{*k} - U \bigr\rVert_{\mathrm{TV}}^2
  \le \frac{1}{4} \sum_{\rho \ne \mathbf{1}} d_\rho \,
  \Tr\bigl(\hat{\mu}(\rho)^k \hat{\mu}(\rho)^{*k}\bigr),
\]
where the sum is over non-trivial irreducible representations $\rho$
of $G$ with dimension $d_\rho$, and $\hat{\mu}(\rho)$ is the Fourier
transform of $\mu$ at $\rho$.
\end{theorem}
